% Research fragment for the parent article. All labels use r8c:.
% Source revision: 73f18ef013d5b74eafc883efb2c15210642df12a.
% New contribution: complete p=8 nonvanishing, leading Hankel factorization.
\section{The complete order-eight reconstruction theorem}
\label{r8c:section}

The incoming reconstruction report \cite[Section 5]{Reconstruction} proves
optimal-prefix reconstruction for $(p,q)=(8,2b)$ when
$b\not\equiv1\pmod {16}$, while explicitly leaving the remaining
residue class open. We prove nonvanishing for that class as well.
The proof identifies the highest Bernoulli-degree part of the exact
four-dimensional determinant as a positive zeta moment determinant.
A uniform comparison handles all $b\ge512$; exact finite arithmetic
then closes the $31$ outstanding smaller cases. The conclusion is an
exact reconstruction theorem and yields ordinary convergent identities
for products of higher polygammas. No arithmetic independence of
special values is used or claimed.

\subsection{Source coordinates and the matrix whose rank is at issue}

Put
\[
 A_s(n)=\zeta(s,n+1)=\zeta(s)-H_n^{(s)},\qquad x_n=n+\tfrac12,
 \qquad D_{b,r}=D_{8,2b}(-2r),
\]
where
\[
 D_{8,2b}(s)=\sum_{n\ge0}A_8(n)A_{2b}(n)x_n^{-s}
\]
is initially defined for $\operatorname{Re}s>-2b-5$ and is continued
by its centered asymptotic expansion. Its nonpositive even arguments
are regular. The inherited inverse Bernoulli coordinates are
\begin{align}
 C_j&=\sum_{r=0}^j\binom{2j+1}{2r}2^{-2(j-r)}D_{b,r},
 \label{r8c:C}\\
 D_{b,r}&=\sum_{j=0}^r\frac{\binom{2r}{2j}}{2j+1}
 B_{2r-2j}(1/2)C_j.
 \label{r8c:D}
\end{align}
For $j\ge b$ these reduce to a rational linear combination of
$Z_k=\zeta(2k)$, $0\le k<b$, where $Z_0=-1/2$.
The exact late block $M_b=(M_{ik})_{0\le i,k<b}$ is
\begin{equation}\label{r8c:M}
 M_{ik}=B_{2i+2k-6}
 \left\{\frac{\binom{2i+2}{8-2k}}{2i+2}
 +\frac{\binom{2b+2i-6}{2b-2k}}{2b+2i-6}\right\}.
\end{equation}
An entry with negative Bernoulli index is zero. A binomial summand is
omitted unless its lower index is at least $2$ and at most its upper
index. This exclusion includes lower index zero. Thus
\[
 (C_b,\ldots,C_{2b-1})^{\mathsf T}
       =M_b(Z_0,\ldots,Z_{b-1})^{\mathsf T}.
\]
More generally, use the same entry formula for every $i\ge0$ and
$0\le k<b$.  The finite Bernoulli-binomial moment formula from
\cite{GeneralHankel} is
\[
 C_j=\sum_{k=0}^{b-1}M_{j-b,k}Z_k\qquad(j\ge b).
\]
Thus these coefficients also specify every later coordinate,
not only the square block used for the rank proof.
The determinant of the first $2b$ moments in the formal coordinate
list $C_0,\ldots,C_{b-1},Z_0,\ldots,Z_{b-1}$ is
$\det M_b/(4b-1)!!$. This is a statement about rational coefficient
rank, before evaluating the coordinates as complex numbers.

\begin{theorem}[All unequal pairs with smaller order eight]
\label{r8c:main}
For every integer $b\ge5$, $\det M_b\ne0$. Consequently, for every
even $q>8$, the moments
\[
 D_{8,q}(0),D_{8,q}(-2),\ldots,D_{8,q}(-2(q-1))
\]
have full formal rank $q$, and every later centered even moment is
an explicit rational linear combination of this prefix.
\end{theorem}

\subsection{An explicitly defined four-dimensional determinant}

We retain the source normalization, to make the comparison of claims
and certificates unambiguous. For $0\le i,\ell\le3$ set
\[
 n_i=b-4+i,\quad N_{i\ell}=2(n_i-\ell)+1,\quad
 R_\ell=2b-6+2\ell,
\]
\[
 F_{i\ell}=\frac1{(2\ell+1)!}
       \prod_{r=0}^{2\ell-1}(2n_i-r)(2n_i+3+r).
\]
An empty product is one. For
\[
 W_{i\ell}(z)=(z+2\ell+1)
                 (2n_i+2\ell+z+1)^{\underline{2i}},
\]
define the Newton coefficients
\begin{equation}\label{r8c:d}
 d_{i\ell h}=\frac1{h!}\sum_{j=0}^{h}
 (-1)^{h-j}\binom hj W_{i\ell}(N_{i\ell}-j).
\end{equation}
The underline means a falling factorial. With independent variables
$U_1,\ldots,U_7$, put
\begin{equation}\label{r8c:smallA}
 \mathcal A_{i\ell}=2F_{i\ell}+
 \frac{R_\ell!}
 {(2n_i+1)(2n_i+2)(2\ell+2)!N_{i\ell}!}
 \sum_{\substack{1\le h\le2i+1\\ h\ \mathrm{odd}}}
 d_{i\ell h}N_{i\ell}^{\underline h}
 U_{4-i+\ell+(h-1)/2}.
\end{equation}
Ratios of factorials with fixed index difference mean their finite
rational products in $b$.

Set
\begin{align}
 D_8(b)={}&256b^8-3584b^7+18496b^6-40544b^5+25984b^4
 \notag\\
 &+18424b^3-5856b^2-8316b-2025,
 \label{r8c:D8}\\
 P(b;U)&=576(2b-5)D_8(b)\frac{\det\mathcal A(b;U)}{\det F(b)}.
 \label{r8c:Pdef}
\end{align}
The source's finite polynomial certificate proves
$P\in\mathbb Z[b,U_1,\ldots,U_7]$, of degree at most four in
$U$, and gives
\begin{align}
 \det F={}&\frac{256}{4725}(b-2)(b-1)(2b-5)(2b-3)
 (4b-11)(4b-9)\notag\\
 &\hspace{8mm}\cdot(4b-7)^2(4b-5)(4b-3)D_8(b).
 \label{r8c:Fdet}
\end{align}
The integer $D_8(b)$ is odd at every integer $b$, so these denominators
are nonzero in the stated range.

Here is its exact connection to the source matrix. Let
\[
 \mu_r=\begin{cases}B_{2r}/(2r)!,&r\ge0,\\0,&r<0,\end{cases}
 \qquad H_b=(\mu_{i+j-3})_{0\le i,j<b},
\]
\[
 R_b=\operatorname{diag}_{0\le i<b}(2b+2i-7)!,\qquad
 E_b=\operatorname{diag}_{0\le k<b}\frac1{(2b-2k)!}.
\]
For every integer $b\ge8$ the inherited reduction is
\begin{equation}\label{r8c:ratio}
 \frac{\det M_b}{\det R_b\det E_b\det H_b}
 =\frac{P(b;2B_{2b-2},\ldots,2B_{2b-14})}
 {576(2b-5)D_8(b)}.
\end{equation}
The matrix $H_b$ is nonsingular. One short proof conjugates it by
$\operatorname{diag}(1,-1,1,-1,\ldots)$ and partitions after the
first two indices. The lower-right block is a positive definite
moment matrix on the points $(2\pi m)^{-2}$, $m\ge1$. The
upper-left block is zero, and the rectangular off-diagonal block
has row rank two: its first two columns form a nonsingular
antitriangular matrix. Its Schur complement is negative definite.
Hence $H_b$ has exactly two negative eigenvalues and is nonsingular.
Thus nonvanishing of the numerator of \eqref{r8c:ratio} proves the
rank assertion.

All formulas in this subsection are inherited definitions or
reductions. The proof below supplies the previously missing global
nonvanishing. In particular, the symbolic certificate is used as a
finite algebraic identity, not as an extrapolation from values of $b$.

\subsection{The highest-degree part is a zeta moment determinant}

Write $P=P_4+P_{\le3}$ according to total degree in $U_1,\ldots,U_7$.
For an integer $b\ge8$, set
\[
 F_0=(2b-14)!,\qquad
 \alpha_j=\frac{(2b-2j)!}{F_0},\qquad V_j=\frac{U_j}{\alpha_j}
       \quad(1\le j\le7),
\]
and define
\begin{align}
 G(b)={}&12288(b-6)^4(b-5)^4(b-4)^4(b-3)^3(b-2)^2(b-1)
 \notag\\
 &\cdot(2b-13)^4(2b-11)^4(2b-9)^4(2b-7)^3
 (2b-5)^3(2b-3)^2(2b-1).
 \label{r8c:G}
\end{align}

\begin{lemma}[Highest-degree Hankel factorization]\label{r8c:hankel}
The exact identity is
\begin{equation}\label{r8c:P4}
 P_4(b;U)=G(b)\det(V_{i+j+1})_{0\le i,j\le3}.
\end{equation}
At the actual doubled Bernoulli values, this becomes
\begin{equation}\label{r8c:P4zeta}
 P_4(b;2B_{2b-2},\ldots,2B_{2b-14})
 =\frac{256G(b)F_0^4}{(2\pi)^{8b-32}}
       \det\bigl(\zeta(2b-14+2i+2j)\bigr)_{0\le i,j\le3}>0.
\end{equation}
\end{lemma}

\begin{proof}
This factorization has a direct structural proof. In
\eqref{r8c:d}, substitute $z=N_{i\ell}-y$. Then
\[
 W_{i\ell}(N_{i\ell}-y)
    =(2n_i+2-y)(4n_i+2-y)^{\underline{2i}},
\]
which is independent of $\ell$. Write its Newton coefficients as
$d_{ih}$. In particular,
\begin{equation}\label{r8c:d1}
 d_{i1}=-(i+1)(4n_i+2)^{\underline{2i}}.
\end{equation}
Indeed the coefficient is the value at $y=1$ minus the value at
$y=0$; using
$(x-1)^{\underline{2i}}=x^{\underline{2i}}(x-2i)/x$
and $4n_i+2=2(2n_i+1)$ gives \eqref{r8c:d1}.

Let $T=\mathcal A-2F$. In a term with $h=2k+1$, the Bernoulli
variable index is $j=4-i+\ell+k$ and
$N_{i\ell}-h=2b-2j$. Thus substitution of $U_j=\alpha_jV_j$
cancels the factorials in \eqref{r8c:smallA}. Define
\[
 \rho_i=\frac1{(2n_i+1)(2n_i+2)},\qquad
 \sigma_\ell=\frac{R_\ell!}{F_0(2\ell+2)!},
\]
and the lower triangular matrix $L$ by
$L_{i,i-k}=d_{i,2k+1}$ for $0\le k\le i$.
Then the exact matrix identity is
\[
 T=\operatorname{diag}(\rho_i)\,L\,
       (V_{4-i+\ell})_{0\le i,\ell\le3}
       \operatorname{diag}(\sigma_\ell).
\]
Reversing its four middle rows has sign $(-1)^6=1$ and gives
the Hankel matrix in \eqref{r8c:P4}. The diagonal entries of $L$
are \eqref{r8c:d1}. Taking determinants, multiplying by
$576(2b-5)D_8(b)/\det F(b)$, and cancelling the finite products
in \eqref{r8c:Fdet} gives exactly \eqref{r8c:G}.
Since the highest-degree part of $\det(2F+T)$ is $\det T$,
this proves \eqref{r8c:P4}. An independent coefficient check in the
accompanying script verifies all sixteen Hankel monomials.

Euler's Bernoulli formula \cite[Eq. 24.8.1]{DLMF} gives
\[
 V_j=(-1)^{b-j+1}\frac{4F_0}{(2\pi)^{2b-2j}}
                     \zeta(2b-2j).
\]
Factoring the row and column signs and powers of $2\pi$ in the
determinant, and then reversing both its rows and its columns,
gives \eqref{r8c:P4zeta}. All signs cancel. Its positivity follows
from the strictly positive moment representation proved next.
\end{proof}

\begin{lemma}[An explicit positive contribution]\label{r8c:positive}
For every real $n>1/2$,
\begin{equation}\label{r8c:zeta-lower}
 \det\bigl(\zeta(2n+2i+2j)\bigr)_{0\le i,j\le3}
 \ge \frac{30625}{48922361856}\,24^{-2n}>0.
\end{equation}
\end{lemma}
\begin{proof}
With $x_m=m^{-2}$, the entries are the moments
$\sum_{m\ge1}x_m^{n+i+j}$. Cauchy--Binet for a finite truncation,
followed by convergence of each entry, yields
\[
 \det\bigl(\zeta(2n+2i+2j)\bigr)
 =\sum_{m_1<m_2<m_3<m_4}
   (m_1m_2m_3m_4)^{-2n}
   \prod_{i<j}(m_i^{-2}-m_j^{-2})^2.
\]
Every summand is positive. Keeping the term $(1,2,3,4)$ proves
\eqref{r8c:zeta-lower}, since
\[
 \prod_{1\le i<j\le4}(i^{-2}-j^{-2})^2
       =\frac{30625}{48922361856}=:V.
\]
No limiting determinant sign is inferred numerically.
\end{proof}

\subsection{Uniform domination of every lower-degree term}

The integer polynomial defined by \eqref{r8c:Pdef} has $94$
monomials when grouped by the seven $U$ variables. Every coefficient
of $P_{\le3}$ has degree at most $15$ in $b$, and its full integer
coefficient norm is exactly
\begin{equation}\label{r8c:K}
 K:=\sum_{\substack{\text{monomials }c b^rU^\nu\text{ of }P\\
                         |\nu|\le3}}|c|
       =1465775272343372.
\end{equation}
This finite integer is obtained by exact coefficient extraction,
independently replayed from the source polynomial certificate.
It is an upper bound on the whole lower-degree contribution, with
no cancellation assumed.

For $b\ge512$, let
\[
 T_b=\frac{8(2b)^{12}(2b-14)!}{(2\pi)^{2b-14}}.
\]
Every actual variable satisfies $|2B_{2b-2j}|\le T_b$:
use $\zeta(2b-2j)<2$, bound the factorial ratio by $(2b)^{12}$,
and use $2\pi>1$ for the remaining nonnegative exponent.
Also $T_b\ge1$. For example $20!>8^{20}$, the ratio for subsequent
even factorials is greater than $8^2$, and $2\pi<8$.
Consequently
\begin{equation}\label{r8c:small-bound}
 |P_{\le3}(b;2B_{2b-2},\ldots,2B_{2b-14})|
 \le Kb^{15}T_b^3.
\end{equation}
For $b\ge14$, every $b-j$ factor in \eqref{r8c:G} is at least
$b/2$, and every $2b-j$ factor is at least $b$. There are $18$
factors of the first kind and $21$ of the second, with multiplicity.
Thus
\[
 G(b)\ge\frac{3}{64}b^{39}.
\]
Combining this with \eqref{r8c:P4zeta} and
\eqref{r8c:zeta-lower} gives
\[
 P_4\ge
 \frac{12Vb^{39}((2b-14)!)^4}{(2\pi)^{8b-32}}\,24^{-2b+14}.
\]
Dividing \eqref{r8c:small-bound} by this expression, first cancelling
the common powers of $2\pi$ and then using $2\pi<8$, proves
\begin{equation}\label{r8c:ratio-bound}
 \frac{|P_{\le3}|}{P_4}
 \le R_b:=C\,\frac{b^{12}192^{2b}}{(2b-14)!},\qquad
 C=\frac{K2^{73}}{3V24^{14}}
   =\frac{211240479050703917773579966480384}{602791875}.
\end{equation}
The cancellation before replacing $2\pi$ is useful: it avoids
artificially weakening the exponential factor.

All that remains here is elementary exact integer arithmetic:
\begin{equation}\label{r8c:base-bound}
 R_{512}<2^{-675}<1,
 \qquad
 \frac{R_{b+1}}{R_b}
 =\left(1+\frac1b\right)^{12}
   \frac{192^2}{(2b-12)(2b-13)}<1\quad(b\ge512).
\end{equation}
The first inequality means the exact integer comparison
\[
 2^{675}\,211240479050703917773579966480384\,
       512^{12}192^{1024}
       <602791875\cdot1010!,
\]
which is verified in the accompanying script. For the second,
$(513/512)^{12}<2$ and $2\cdot192^2<1011\cdot1012$ suffice.
It follows that $P=P_4+P_{\le3}>0$ for every $b\ge512$.
The threshold is chosen for a transparent uniform proof; no minimal
threshold is asserted.

\subsection{Finite completion of the omitted residue class}

For completeness, the inherited congruence proof gives, for $b\ge11$,
\[
 P(b;2B_{2b-2},\ldots,2B_{2b-14})
       \equiv4(b-1)(4b+15)\pmod {64}
       \quad\text{in }\mathbb Z_{(2)}.
\]
It uses $2B_{2m}\equiv1+4m+16m^2\pmod {64}$ for $m\ge4$,
together with its symbolic determinant certificate. Thus $P\ne0$
when $b\not\equiv1\pmod {16}$. The source checks $5\le b\le10$
directly. Its congruence assertion does not supply the omitted
residue class, and no extrapolation of its valuation formula is
used here.

Below $512$ the only remaining indices are $b=16k+1$,
$1\le k\le31$. The following table records the exact valuation
of $P$ at every such index.
\begin{center}
\begin{tabular}{r r@{\qquad}r r@{\qquad}r r@{\qquad}r r}
\hline
$b$&$v_2(P)$&$b$&$v_2(P)$&$b$&$v_2(P)$&$b$&$v_2(P)$\\
\hline
17&6&33&7&49&6&65&10\\
81&6&97&7&113&6&129&8\\
145&6&161&7&177&6&193&9\\
209&6&225&7&241&6&257&8\\
273&6&289&7&305&6&321&11\\
337&6&353&7&369&6&385&8\\
401&6&417&7&433&6&449&9\\
465&6&481&7&497&6&&\\
\hline
\end{tabular}
\end{center}
Every entry is certified by a nonzero residue modulo $2^{32}$;
all $31$ residues are included in the machine-readable record.
They can be obtained either from exact rational Bernoulli numbers
or entirely in modular arithmetic using
\[
 (2m+1)\,(2B_{2m})
 =2m-1-\sum_{j=1}^{m-1}\binom{2m+1}{2j}\,(2B_{2j}).
\]
The divisor $2m+1$ is odd, so the recurrence is valid modulo
$2^{32}$. Two independent implementations agree on every residue.
In particular no finite numerical approximation is treated as a
nonvanishing certificate. As an explicit first new value,
\[
 \frac{\det M_{17}}{\det R_{17}\det E_{17}\det H_{17}}
   =-\frac{5564518155020169678178976040991}{8937041637533625}.
\]
This settles the previously omitted pair $(p,q)=(8,34)$.
The inherited cases, this finite table, and
\eqref{r8c:base-bound} cover every $b\ge5$, proving
Theorem~\ref{r8c:main}.

\subsection{An all-order family of ordinary polygamma identities}

The rank theorem supplies explicit identities as well as a decision
about a matrix. For $b\ge5$ and $0\le k<b$, define
\begin{align}
 w_{k,r}^{(b)}&=
 \sum_{t=0}^{b-1}(M_b^{-1})_{kt}
 \binom{2(b+t)+1}{2r}2^{-2(b+t-r)},
 \label{r8c:w}\\
 Q_{b,k}(x)&=\sum_{r=0}^{2b-1}w_{k,r}^{(b)}x^{2r},
 \label{r8c:Q}
\end{align}
where the summand is zero for $r>b+t$. These are finite rational
formulas for every coefficient. Put
\[
 \eta_{s,0}=\frac1{s-1},\qquad
 \eta_{s,j}=\frac{(s)_{2j-1}B_{2j}(1/2)}{(2j)!}\quad(j\ge1),
 \qquad
 c_{b,\ell}=\sum_{j=0}^{\ell}\eta_{8,j}\eta_{2b,\ell-j}.
\]
Let $S_{b,k}(x)$ be the polynomial formed by retaining the
nonnegative powers in the formal product
\[
 Q_{b,k}(x)x^{-2b-6}\sum_{\ell\ge0}c_{b,\ell}x^{-2\ell}.
\]
Only finitely many terms enter this polynomial, and its degree is
at most $2b-8$. Then
\begin{equation}\label{r8c:ordinary}
 \boxed{\sum_{n=0}^{\infty}
 \left\{Q_{b,k}(x_n)
 \frac{\psi^{(7)}(n+1)\psi^{(2b-1)}(n+1)}{7!\,(2b-1)!}
       -S_{b,k}(x_n)\right\}=\zeta(2k).}
\end{equation}
This is an ordinary absolutely convergent series. Indeed
\eqref{r8c:w} and \eqref{r8c:C} identify its continued moment
with the $k$th row of $M_b^{-1}M_b$, giving $\zeta(2k)$.
The centered Hurwitz expansion is
\[
 A_8(n)A_{2b}(n)\sim x_n^{-2b-6}
                      \sum_{\ell\ge0}c_{b,\ell}x_n^{-2\ell}.
\]
The polynomial $S_{b,k}$ removes every nondecaying term, leaving
$O(x_n^{-2})$. Each removed nonnegative even power has centered
continued sum $\zeta(-2j,1/2)=0$, so the ordinary convergent sum
has precisely the same value. Thus no unspecified summation of a
divergent series remains in \eqref{r8c:ordinary}.

To reconstruct a later moment, first use \eqref{r8c:C} on the
initial prefix, invert $M_b$, form every later $C_j$ from the
source Bernoulli-binomial formula, and finally use \eqref{r8c:D}.
Every operation is finite rational arithmetic. This completes the
order-eight family. The still-uncovered pairs have smaller even order
at least ten and order gap at least six; their optimal-prefix
conjecture remains separate.
